\chapter{Grading and the Algebra of Novelty}\label{ch:grading}

\epigraph{\itshape ``More is different.''}{P. W. Anderson \cite{Anderson1972}}

The previous chapters built the static apparatus of an idea algebra: a
monoid of compositions $\IdMon$, a resonance pairing $\rho$, an Aufhebung
decomposition $a\compose b = \preserve(a,b) + \negate(a,b) +
\elevate(a,b)$, and an emergence cocycle $\emerg$. These data describe
\emph{what} ideas are and how they combine, but not \emph{how far} a
given combination has travelled from its constituents. Travel is the
business of the grading. In this chapter we promote the informal
practice of speaking about ``higher-order'' ideas, ``second-generation''
theories, and ``levels of abstraction'' to a precise piece of structure:
a monoid homomorphism
\[
  \deg : \Ideas \longrightarrow G
\]
into a totally ordered abelian group $G$, compatible with every other
operation of the algebra. We then read the predicate of \emph{novelty}
off the grading: $a$ is novel relative to $b$ iff $\deg(a) > \deg(b)$.
The chapter culminates in three theorems --- the
Reassociation Theorem, the Cumulative-Novelty Theorem, and the
Transmission-Preservation Theorem --- and in a spectral decomposition of
$\Ideas$ by degree that we will use throughout the rest of the book. We
close, as always, with a zoo of historical positions whose intuitions
the formal grading recovers.

\section{Why grade? From ``level of an idea'' to a homomorphism}
\label{sec:why-grade}

\subsection*{The pre-formal practice}

Almost every classical theory of ideas talks about ranks. Plato has
levels of being; Aquinas distinguishes \emph{species intelligibilis} of
graded abstraction; Lovejoy writes of unit-ideas combining into
elaborate doctrines \cite{Lovejoy1936}; Lakatos grades research
programmes by problem-shift sophistication \cite{Lakatos1970}; Romer
counts recombinations \cite{Romer1990}; Henrich grades cultural
inheritance by generation \cite{Henrich2015}. Each tradition assigns,
explicitly or tacitly, a numerical or ordinal index to ideas, and each
expects this index to behave additively under combination: the level of
a synthesis is at least the sum of the levels of its parts. The thesis
of this chapter is that all of these intuitions are instances of a
\emph{single} structural axiom --- the existence of a homomorphism
$\deg$ into an ordered abelian group --- and that the apparently
disparate ``zoo'' of grading practices is forced by that axiom together
with the compatibility laws below.

What is gained by formalisation? Three things. First, a clean criterion
for when a putative grading is \emph{coherent} with the rest of the
algebra: it must respect $\compose$, the Aufhebung pieces $\preserve,
\negate, \elevate$, and the cocycle $\emerg$. Many naive gradings ---
``count the words of the idea'', ``count the citations'', ``count the
neurons that fire'' --- fail this criterion and so cannot be the
$\deg$ of an idea algebra. Second, an automatic notion of novelty: once
$\deg$ is in place, the predicate ``$a$ is novel relative to $b$'' has
a forced definition, and the corresponding theorems about novelty
(\cref{sec:novelty,sec:headline}) are theorems about $G$, not free
choices. Third, a spectral decomposition of $\Ideas$ by degree
(\cref{sec:spectral}), which makes the Anderson slogan ``more is
different'' \cite{Anderson1972} into a precise statement about a graded
algebra and gives Bedau's notion of weak emergence \cite{Bedau1997} the
form $\emerg \in \ker \deg$.

\subsection*{Why a homomorphism, rather than a mere function?}

One might object that we could in principle define $\deg$ as any
function $\Ideas \to G$ without imposing the homomorphism axiom (G1).
The answer is that without (G1), the spectral decomposition of
\cref{prop:spectral} loses its multiplicative structure, the
Reassociation Theorem (\cref{thm:reassociation}) fails to give a
coherent assignment of degree to traditions, and the very notion of
``cumulative novelty'' becomes uninterpretable. The homomorphism
axiom is what makes the grading \emph{compositional}: the degree of a
composite can be predicted from the degrees of its parts (modulo a
single cocycle correction), and this is the formal counterpart of the
intuition that ``higher-level'' ideas are built out of ``lower-level''
ones in a way that respects level. A grading without (G1) is, in
practice, a typology --- a partition of $\Ideas$ into kinds --- not a
yardstick.

A subtler point is that (G1) does not force $\deg$ to be a
\emph{strict} homomorphism. The cocycle term $\omega^G$ allows for
``emergence corrections'' that prevent the degree of a composite from
being the simple sum of constituent degrees. This is a feature, not a
bug: real intellectual history is full of cases where a synthesis
acquires a degree that is neither the sum nor the maximum of its
inputs, and the cocycle is the formal device that records exactly how
much this departure costs. The Grading Theorem
(\cref{thm:grading}) shows that the cocycle is uniquely determined by
$\deg$, so the formalism does not multiply parameters: choosing
$\deg$ choices $\omega^G$.

\subsection*{What the total order on $G$ buys}

The decision to take $G$ totally ordered, rather than merely partially
ordered or lattice-valued, has substantive content beyond
mathematical convenience. A partial order on $G$ would permit two
ideas $a, b$ with incomparable degrees, leaving open the question
``which is more novel?'' --- a question the chapter is meant to
answer definitively. Total order rules this out. The cost is that
some natural multidimensional measures of complexity (e.g.\ a
combination of structural depth and information content) cannot be
the codomain of a single grading; they must be projected to a
totally ordered scale, with attendant loss of information.

This loss is not always undesirable. The historical practice of
ranking ideas (``higher'', ``deeper'', ``more general'') is itself
totally ordered, and a formalisation that respects this practice is
empirically apt for the historical material the zoo covers. When
multidimensional gradings are needed --- for instance, in the
analysis of conceptual spaces \cite{Gardenfors2000} --- one can
extend the present formalism by taking $G$ to be a lexicographically
ordered product, recovering the desired refinement while preserving
total order.

\section{Graded idea algebras}\label{sec:graded-axioms}

We fix throughout this chapter a totally ordered abelian group
$(G,+,0,\le)$. The leading examples are $G = \ZZ$ (counting generations
or recombinations), $G = \QQ$ (allowing fractional contributions, e.g.\ a
co-authored idea), and $G = \RR$ (admitting continuous valuations such
as information content).

\begin{definition}[Graded idea algebra]\label{def:graded-idea-alg}
A \emph{graded idea algebra} is a tuple
$(\Ideas,\compose,\unit,\rho,\preserve,\negate,\elevate,\emerg,\deg)$
where $(\Ideas,\compose,\unit,\rho,\preserve,\negate,\elevate,\emerg)$
is an idea algebra in the sense of Theorems 1--8, $G$ is a totally
ordered abelian group, and $\deg : \Ideas \to G$ is a function
satisfying the following compatibility axioms.
\begin{enumerate}
\item[\textbf{(G1)}] \emph{Monoid homomorphism}:
  $\deg(\unit) = 0$ and
  $\deg(a\compose b) = \deg(a) + \deg(b) + \omega^G(a,b)$,
  where $\omega^G : \Ideas\times\Ideas \to G$ is the image of $\emerg$
  along the canonical map $A \to G$ when $A=G$ (and is $0$ otherwise).
\item[\textbf{(G2)}] \emph{Resonance compatibility}: if $\rho(a,a') \neq
  0$ then $\deg(a) = \deg(a')$. Resonance is a degree-preserving
  pairing.
\item[\textbf{(G3)}] \emph{Aufhebung compatibility}:
  $\deg(\preserve(a,b)) \le \min(\deg(a),\deg(b))$,
  $\deg(\negate(a,b))   \le \max(\deg(a),\deg(b))$,
  $\deg(\elevate(a,b))  \ge \max(\deg(a),\deg(b))$.
  Preservation does not raise degree, negation does not exceed the
  maximum, and elevation/synthesis only ever raises it.
\item[\textbf{(G4)}] \emph{Cocycle compatibility}: $\omega^G$ is itself a
  $G$-valued normalized $2$-cocycle, i.e.\ $\dcoc \omega^G = 0$ and
  $\omega^G(\unit,a) = \omega^G(a,\unit) = 0$.
\item[\textbf{(G5)}] \emph{Identity-non-novelty}: $\deg(\unit)$ is the
  unique minimum of $\deg$ on the set
  $\{a\in\Ideas : \exists b,\ a = \unit\compose b \text{ or } a =
  b\compose \unit\}$. Equivalently, the silent idea has the lowest
  possible degree among ideas reachable from itself by composition.
\end{enumerate}
\end{definition}

\begin{remark}[The $\omega^G$ correction]
When $A=G$, axiom (G1) says $\deg$ is a $1$-cochain whose coboundary is
$\dcoc \deg(a,b) = \deg(a) + \deg(b) - \deg(a\compose b)$, and (G1)
identifies this coboundary with $-\omega^G$. Thus a grading is a
$1$-cochain that \emph{trivializes} the negative of the cocycle. When
$\emerg$ is cohomologically trivial (e.g.\ $\emerg=\dcoc v$ for a
valuation $v$), one may take $\deg = v$ and obtain a strict
homomorphism. When $\emerg$ is non-trivial, no strict homomorphism
exists and $\omega^G$ is forced to be the actual coboundary of $\deg$.
This is the formal sense in which ``emergence is the obstruction to
gradeability'' alluded to by Bedau \cite{Bedau1997}.
\end{remark}

\begin{remark}[Why $G$ is totally ordered]
Lattice-valued or partially ordered codomains permit ``incomparable''
ideas, but rule out a sharp definition of novelty: there is no fact of
the matter whether $a$ is more novel than $b$ when $\deg(a)$ and
$\deg(b)$ are incomparable. The whole point of the grading axiom is to
import a linear yardstick. The total order on $G$ is therefore not a
mathematical convenience but a substantive commitment of the theory.
\end{remark}

\section{The grading theorem}\label{sec:grading-theorem}

We restate Theorem~9 in the present notation.

\begin{theorem}[Theorem 9; Grading Theorem]\label{thm:grading}
Let $(\Ideas,\compose,\unit,\rho,\preserve,\negate,\elevate,\emerg)$
be an idea algebra with abelian coefficient group $A$. Let $G$ be a
totally ordered abelian group and let $\deg_0 : \Ideas \to G$ be a
function satisfying $\deg_0(\unit) = 0$ together with axioms
\textup{(G2)--(G5)} of \cref{def:graded-idea-alg}. Define
\[
  \omega^G(a,b) \defeq \deg_0(a\compose b) - \deg_0(a) - \deg_0(b).
\]
Then $\omega^G$ is a normalized $G$-valued $2$-cocycle, the tuple
$(\Ideas,\compose,\unit,\rho,\preserve,\negate,\elevate,\emerg,\deg_0)$
is a graded idea algebra, and any other extension $\deg_1$ of
$\deg_0|_{\{\unit\}}$ satisfying the same axioms differs from
$\deg_0$ by an additive constant on each connected component of
$\Ideas$ in the composition graph.
\end{theorem}

\begin{proof}
We verify the four conclusions in turn.

\emph{(i) $\omega^G$ is a $2$-cocycle.} Compute, using associativity of
$\compose$,
\begin{align*}
  \omega^G(a,b) + \omega^G(a\compose b, c)
  &= \bigl[\deg_0(a\compose b) - \deg_0(a) - \deg_0(b)\bigr] \\
  &\quad{} + \bigl[\deg_0((a\compose b)\compose c) - \deg_0(a\compose b)
            - \deg_0(c)\bigr] \\
  &= \deg_0((a\compose b)\compose c) - \deg_0(a) - \deg_0(b) - \deg_0(c).
\end{align*}
By associativity, $(a\compose b)\compose c = a\compose(b\compose c)$, so
the same expression equals $\omega^G(b,c) + \omega^G(a,b\compose c)$.
Subtracting yields the cocycle identity $\dcoc \omega^G = 0$.
Normalisation $\omega^G(\unit,a) = \omega^G(a,\unit) = 0$ follows
because $\deg_0(\unit) = 0$ and $\unit\compose a = a\compose \unit = a$
in the underlying monoid.

\emph{(ii) Axiom (G1).} By the definition of $\omega^G$, we have
$\deg_0(a\compose b) = \deg_0(a) + \deg_0(b) + \omega^G(a,b)$, which is
exactly (G1).

\emph{(iii) The remaining axioms.} Axioms (G2)--(G5) hold by
hypothesis on $\deg_0$.

\emph{(iv) Uniqueness up to constants.} Suppose $\deg_1$ is another
function satisfying (G1)--(G5) with the same induced cocycle (the
hypothesis ``with the same axioms'' fixes the cocycle as
$\omega^G_1(a,b) \defeq \deg_1(a\compose b) - \deg_1(a) - \deg_1(b)$).
Define $f \defeq \deg_1 - \deg_0$. Then for all $a,b$,
\[
  f(a\compose b) = f(a) + f(b) + \bigl(\omega^G_1(a,b) - \omega^G(a,b)\bigr).
\]
If $\omega^G_1 = \omega^G$, then $f$ is a monoid homomorphism into
$G$. The composition graph of $\Ideas$ has the same connected components
as the underlying monoid; on each component, $f$ is determined by its
value on a single representative. By (G5) we have $f(\unit) = 0$ on the
identity component, and by additivity $f$ is constant on each component
of the composition graph rooted at $\unit$. Thus $\deg_1$ and $\deg_0$
agree up to a constant on each connected component, as claimed.
\end{proof}

\begin{remark}[The role of total order]
In the proof, totality of $\le$ on $G$ is used only in the
\emph{interpretation} of (G3), which involves $\min$ and $\max$. The
arithmetic identities --- cocycle, normalisation, uniqueness --- are
purely additive and work for any abelian $G$. The total order is what
makes ``the degree drops'', ``the degree rises'' meaningful, and what
will let us speak of novelty in \cref{sec:novelty}.
\end{remark}

\begin{corollary}[Trivial emergence and homomorphic gradings]
\label{cor:trivial-grading}
If $\emerg$ is the coboundary of a $G$-valued $1$-cochain $v$, i.e.\
$\emerg(a,b) = v(a) + v(b) - v(a\compose b)$, then $\deg \defeq v$
satisfies $\omega^G \equiv 0$, and (G1) reduces to the strict
homomorphism law $\deg(a\compose b) = \deg(a) + \deg(b)$.
\end{corollary}

\begin{proof}
Substitute $v$ for $\deg_0$ in the definition of $\omega^G$ and observe
that the result is identically zero.
\end{proof}

\section{Novelty}\label{sec:novelty}

The grading lets us define novelty without further appeal to the
content of any specific idea.

\begin{definition}[Novelty predicates]\label{def:novelty}
Let $(\Ideas,\dots,\deg)$ be a graded idea algebra and let $a,b\in
\Ideas$.
\begin{itemize}
\item $a$ is \emph{novel relative to} $b$, written $a \succ b$, iff
  $\deg(a) > \deg(b)$.
\item $a$ is \emph{strictly novel} iff $\deg(a) > 0$, i.e.\ $a$
  is novel relative to the silent idea $\unit$.
\item Given a finite list $L = (b_1,\dots,b_n)$, the idea $a$ is
  \emph{novel-everywhere on $L$} iff $a \succ b_i$ for all
  $i\in\{1,\dots,n\}$.
\item A list $L = (a_1,\dots,a_n)$ is a \emph{chain-novel} or
  \emph{ascending lineage} iff $a_1 \prec a_2 \prec \cdots \prec a_n$.
\end{itemize}
\end{definition}

The Lean development of Theorems~10--12 records a closely related
predicate $\mathsf{IsNovel}(c,a) \defeq c \neq a$, together with a
strengthened version $\mathsf{EmergentNovel}(a,b) \defeq a\compose b
\notin \{a,b\}$ that asks the composition to differ from \emph{both}
ingredients. Under any grading respecting axiom (G3), these inequalities
are detected by the degree.

\begin{lemma}[Emergent novelty implies degree increase]
\label{lem:emergent-implies-degree}
Suppose $(\Ideas,\dots,\deg)$ is a graded idea algebra and $a,b\in
\Ideas$ satisfy $a\compose b \neq a$ and $a\compose b \neq b$
($\mathsf{EmergentNovel}$). Then either
\[
  \deg(a\compose b) > \max(\deg(a),\deg(b))
  \quad\text{or}\quad
  \deg(a\compose b) < \min(\deg(a),\deg(b)),
\]
unless $\deg$ is constant on the connected component of $a$ and $b$ in
the composition graph.
\end{lemma}

\begin{proof}
By axiom (G1), $\deg(a\compose b) = \deg(a) + \deg(b) + \omega^G(a,b)$.
If $\deg(a\compose b) = \max(\deg(a),\deg(b))$, then in particular,
say WLOG $\deg(a) \ge \deg(b)$, we have $\deg(b) + \omega^G(a,b) = 0$.
By (G3) applied to $\elevate(a,b)$, this forces
$\elevate(a,b)$ to lie at degree $\le \deg(a)$, and the Aufhebung
decomposition $a\compose b = \preserve(a,b) + \negate(a,b) +
\elevate(a,b)$ collapses --- via uniqueness mod $\ker\rho$
(Theorem~3) --- to either $a$ or $b$. This contradicts
$\mathsf{EmergentNovel}(a,b)$, unless every term has the same degree
(the constant case).
\end{proof}

\begin{lemma}[Identity cannot originate novelty]
\label{lem:identity-not-novel}
For every $b\in\Ideas$, $\unit\compose b = b$ and $b\compose\unit = b$,
so neither $(\unit, b)$ nor $(b, \unit)$ is emergently novel.
\end{lemma}

\begin{proof}
By the monoid laws, $\unit\compose b = b = b\compose\unit$. In
particular $\unit\compose b \in \{\unit, b\}$, so the
$\mathsf{EmergentNovel}$ condition fails, and by axiom (G5), $\deg(b) =
\deg(\unit) + \deg(b) + \omega^G(\unit,b) = \deg(b)$, consistent with
$\omega^G(\unit,b) = 0$.
\end{proof}

\section{The headline reassociation and novelty theorems}
\label{sec:headline}

We now translate the principal results of Theorems~10--12. We continue
to write $\mathsf{tradition}(L)$ for the iterated composition
$a_1\compose a_2\compose \cdots\compose a_n$ of a finite list $L =
(a_1,\dots,a_n)$, with $\mathsf{tradition}(\emptyset) = \unit$ and
$\mathsf{tradition}(L_1 \mathbin{\#} L_2) = \mathsf{tradition}(L_1)
\compose \mathsf{tradition}(L_2)$, where $\mathbin{\#}$ denotes list
concatenation.

\begin{theorem}[Reassociation Theorem; Theorem~10.1]
\label{thm:reassociation}
For any three lists $L_1, L_2, L_3$ of ideas,
\begin{enumerate}
\item[(i)] $\mathsf{tradition}(L_1 \mathbin{\#} L_2) =
  \mathsf{tradition}(L_1) \compose \mathsf{tradition}(L_2)$, and
\item[(ii)] $\mathsf{tradition}((L_1 \mathbin{\#} L_2) \mathbin{\#} L_3)
  = \mathsf{tradition}(L_1 \mathbin{\#} (L_2 \mathbin{\#} L_3))$.
\end{enumerate}
Consequently, $\deg(\mathsf{tradition}(L)) = \sum_{a\in L} \deg(a) +
\sum_{i<n} \omega^G(\mathsf{tradition}(L_{<i}), a_i)$ for $L =
(a_1,\dots,a_n)$.
\end{theorem}

\begin{proof}
Both halves are by induction on $|L_1|$. For (i), if $L_1 = \emptyset$,
both sides equal $\mathsf{tradition}(L_2)$; for $L_1 = a :: L_1'$,
\begin{align*}
  \mathsf{tradition}((a :: L_1') \mathbin{\#} L_2)
  &= a \compose \mathsf{tradition}(L_1' \mathbin{\#} L_2) \\
  &= a \compose (\mathsf{tradition}(L_1') \compose
       \mathsf{tradition}(L_2))
   \quad\text{(induction)}\\
  &= (a \compose \mathsf{tradition}(L_1')) \compose
       \mathsf{tradition}(L_2)
   \quad\text{(associativity)}\\
  &= \mathsf{tradition}(a :: L_1') \compose \mathsf{tradition}(L_2).
\end{align*}
Part (ii) follows from (i) applied twice and associativity of
$\compose$. The degree formula is then immediate by repeated
application of (G1).
\end{proof}

\begin{theorem}[Silent-Agent Invariance; Theorem~10.2]
\label{thm:silent-invariance}
Let $L = (a_1,\dots,a_n)$ be a list and let $L'$ result from $L$ by
inserting any number of copies of $\unit$ at any positions. Then
$\mathsf{tradition}(L) = \mathsf{tradition}(L')$, and in particular
$\deg(\mathsf{tradition}(L)) = \deg(\mathsf{tradition}(L'))$.
\end{theorem}

\begin{proof}
By induction on the number of insertions. Inserting $\unit$ at position
$k$ replaces a substring $a_k\compose \cdots$ by $\unit\compose
a_k\compose \cdots = a_k\compose \cdots$ via (the monoid law). The
degree statement follows by Theorem~9 (G1) and $\omega^G(\unit,b) = 0$.
\end{proof}

\begin{corollary}[Social congruence]\label{cor:social-cong}
The relation $L_1 \sim L_2 :\Leftrightarrow \mathsf{tradition}(L_1) =
\mathsf{tradition}(L_2)$ is an equivalence on finite lists of ideas,
respected by concatenation on either side and by silent insertion. In
particular, $\sim$ descends to an equivalence on degrees:
$\deg\circ\mathsf{tradition} : \mathrm{List}(\Ideas)/{\sim} \to G$ is
well-defined.
\end{corollary}

\begin{proof}
Reflexivity, symmetry and transitivity are immediate. Compatibility
with concatenation is \cref{thm:reassociation}(i). Compatibility with
silent insertion is \cref{thm:silent-invariance}.
\end{proof}

\begin{theorem}[Cumulative Novelty; Theorem~12.2]
\label{thm:cumulative-novelty}
Let $L = (a_1,\dots,a_n)$ be a list of ideas. Say $L$ is
\emph{lineage-novel} iff $\mathsf{tradition}(L) \neq a_i$ for every $i$.
Then:
\begin{enumerate}
\item[(i)] If $L$ is lineage-novel, then $\mathsf{tradition}(L) \notin
  L$, i.e.\ no single member produced the tradition.
\item[(ii)] No singleton $L = (a)$ is lineage-novel.
\item[(iii)] No nonempty list of silent agents $L = (\unit,\dots,\unit)$
  is lineage-novel.
\item[(iv)] The empty list is vacuously lineage-novel.
\end{enumerate}
\end{theorem}

\begin{proof}
(i) is the contrapositive of the assumption: if $\mathsf{tradition}(L) =
a_i$ for some $i$, then $\mathsf{tradition}(L) \in L$. (ii) is because
$\mathsf{tradition}((a)) = \unit\compose a = a$, so the unique member
\emph{is} the tradition. (iii): for $L = (\unit,\dots,\unit)$,
$\mathsf{tradition}(L) = \unit$ by repeated application of $\unit\compose
\unit = \unit$, and $\unit\in L$. (iv) holds vacuously: there is no $i$
to violate the condition.
\end{proof}

\begin{theorem}[Transmission-Preservation of Novelty; Theorem~12.3]
\label{thm:transmission-preservation}
Let $L_1, L_2$ be lists of ideas with traditions $t_1, t_2$. The
following hold.
\begin{enumerate}
\item[(i)] If $t_1 \compose t_2 \notin \{t_1, t_2\}$ (i.e.\ $t_1$ and
  $t_2$ are emergently novel as a pair), then
  $\mathsf{tradition}(L_1\mathbin{\#}L_2) \neq t_1$ and
  $\mathsf{tradition}(L_1\mathbin{\#}L_2) \neq t_2$.
\item[(ii)] If the vertical-transmission map $T_{L} : x \mapsto
  \mathsf{tradition}(L)\compose x$ is faithful (i.e.\ injective), then
  $\mathsf{tradition}(L)\compose x \neq x$ implies
  $\mathsf{tradition}(L)\neq \unit$. In particular, no faithful
  transmission can be right-novel relative to a non-trivial input
  unless the tradition is non-trivial.
\item[(iii)] If $L_1 \sim L_2$ (same tradition), then for every $c$,
  $t_1$ is emergently novel with $c$ iff $t_2$ is.
\end{enumerate}
\end{theorem}

\begin{proof}
(i) By \cref{thm:reassociation}(i), $\mathsf{tradition}(L_1\mathbin{\#}
L_2) = t_1\compose t_2$. The hypothesis says $t_1\compose t_2 \notin
\{t_1,t_2\}$, so directly
$\mathsf{tradition}(L_1\mathbin{\#}L_2) \neq t_1$ and $\neq t_2$.

(ii) Assume $T_L$ is faithful. The Lean
lemma $T_L = \mathrm{leftMul}(\mathsf{tradition}(L))$ (Corollary~11.4)
shows that $T_L(x) = \mathsf{tradition}(L)\compose x$. Faithfulness of
$T_L$ means $T_L(x) = T_L(y) \Rightarrow x=y$. Suppose
$\mathsf{tradition}(L) = \unit$. Then $T_L = \mathrm{id}$, so for every
$x$, $T_L(x) = x$, contradicting $\mathsf{tradition}(L)\compose x \neq
x$. Hence $\mathsf{tradition}(L)\neq \unit$, as claimed.

(iii) By \cref{cor:social-cong}, $L_1\sim L_2$ implies $t_1 = t_2$.
Then for any $c$, $t_1\compose c = t_2\compose c$ and $t_1 = t_2$,
$c=c$, so the predicate $t\compose c \notin \{t,c\}$ has the same truth
value for $t = t_1$ and $t=t_2$.
\end{proof}

\begin{corollary}[Faithful transmission preserves the degree of the input]
\label{cor:faithful-degree}
If $T_L$ is faithful and $\mathsf{tradition}(L) = \unit$, then for every
$x$, $\deg(T_L(x)) = \deg(x)$.
\end{corollary}

\begin{proof}
$T_L(x) = \unit\compose x = x$.
\end{proof}

\subsection*{Refinements: chain-novelty, pointwise-novelty, and the
covering relation}

The novelty predicates of \cref{def:novelty} admit several
sharpenings that are useful in applications. We collect them here.

\begin{definition}[Pointwise novelty]\label{def:pointwise-novelty}
Let $L = (a_1,\dots,a_n)$ be a finite list. A point $a\in\Ideas$ is
\emph{pointwise-novel against $L$} iff $a \succ a_i$ for some
$i\in\{1,\dots,n\}$. It is \emph{novel-everywhere on $L$} (cf.\
\cref{def:novelty}) iff $a\succ a_i$ for all $i$. Pointwise novelty
is the existential weakening; novel-everywhere is the universal
strengthening.
\end{definition}

\begin{lemma}[Quantifier monotonicity]\label{lem:quant-mono}
If $a$ is novel-everywhere on $L$, then $a$ is pointwise-novel
against $L$. Conversely, if $|L|=1$, the two notions coincide.
\end{lemma}

\begin{proof}
The first claim is logically immediate: $\forall i$ implies
$\exists i$ when the index set is non-empty. The second is the
observation that for a singleton $L = (b)$, both predicates reduce to
$a\succ b$.
\end{proof}

\begin{definition}[Covering relation in $G$]\label{def:cover}
The \emph{covering relation} on the image of $\deg$ is
$g \prec\!\!\cdot\, h$ iff $g < h$ and there is no $g'$ with $g < g' <
h$ in the image of $\deg$. An idea $a$ is a \emph{minimal-novelty
extension} of $b$ iff $\deg(b) \prec\!\!\cdot\, \deg(a)$.
\end{definition}

\begin{proposition}[Existence of minimal-novelty extensions]
\label{prop:min-extension}
Suppose the image of $\deg$ in $G$ is order-isomorphic to a sub-monoid
of $\NN$ (so the covering relation is well-defined and
nontrivial). Then for every $b\in\Ideas$ such that $\Ideas_{>\deg(b)}
\neq\emptyset$, there exists a minimal-novelty extension $a$ of $b$.
\end{proposition}

\begin{proof}
By assumption, the image of $\deg$ in $\NN$ is well-ordered, so the
set $\{g\in\mathrm{im}(\deg) : g > \deg(b)\}$ has a minimum, say $g^*$.
Any $a\in\Ideas_{g^*}$ satisfies $\deg(b)\prec\!\!\cdot\,\deg(a)$.
\end{proof}

\begin{remark}[Boden minimality vs.\ chain-novelty]
The covering relation captures a weaker notion than chain-novelty: an
ascending lineage $a_1\prec\cdots\prec a_n$ may skip strata. Boden's
exploratory creativity \cite{Boden2004} corresponds to ascending one
covering step at a time, while transformational creativity corresponds
to skipping strata --- the latter being possible only when the cocycle
$\omega^G$ has a sufficiently large positive value.
\end{remark}

\subsection*{Application: a cumulative-novelty zero-one law}

\begin{theorem}[Zero-one law for cumulative novelty]
\label{thm:zero-one}
Let $L = (a_1, a_2, \dots)$ be an infinite lineage in a graded idea
algebra with $G = \ZZ$ and $\omega^G \equiv 0$. Suppose each $a_i$ is
chosen independently from a stationary distribution on $\Ideas$ with
$\mathbb{E}[\deg(a_i)] = \mu$. Then either $\mu > 0$ and
$\deg(\mathsf{tradition}((a_1,\dots,a_n))) \to +\infty$ almost surely,
or $\mu \le 0$ and the lineage is eventually periodic mod equivalence.
\end{theorem}

\begin{proof}[Proof sketch]
By (G1) with $\omega^G \equiv 0$, $\deg(\mathsf{tradition}(L_n)) =
\sum_{i\le n}\deg(a_i)$. By the strong law of large numbers, this
divided by $n$ tends to $\mu$ almost surely. For $\mu > 0$ the sum
diverges to $+\infty$; for $\mu < 0$ to $-\infty$, contradicting
$\deg \ge 0$ on a generating set, hence the lineage must enter a
recurrent class. The boundary case $\mu = 0$ requires more care
(reduces to a recurrent random walk on $\ZZ$).
\end{proof}

This is a quantitative statement of the Henrich ratchet
\cite{Henrich2015}: cumulative culture requires a positive expected
contribution per generation, and any society with $\mu \le 0$ will
fail to accumulate degree, regardless of the size of its repertoire.

\begin{corollary}[Composite-tradition emergence; Theorem~12.4]
\label{cor:composite-emergence}
Under the hypothesis of
\cref{thm:transmission-preservation}(i),
$\deg(\mathsf{tradition}(L_1\mathbin{\#}L_2)) > \max(\deg t_1,\deg t_2)$
provided the algebra has Aufhebung-positive emergence (i.e.\
$\elevate(t_1,t_2)$ has degree strictly greater than
$\max(\deg t_1,\deg t_2)$).
\end{corollary}

\begin{proof}
Combine \cref{lem:emergent-implies-degree} with axiom (G3).
\end{proof}

\begin{corollary}[Closed-set non-emergence; Theorem~12.5]
\label{cor:closed-no-emergence}
If $S\subseteq \Ideas$ is closed under $\compose$ (i.e.\ $a,b\in S
\Rightarrow a\compose b \in S$), then no $\compose$-product of members
of $S$ is novel relative to $S$ in the set-theoretic sense; equivalently,
the restriction of $\deg$ to $S$ takes values in a sub-monoid of $G$.
\end{corollary}

\begin{proof}
Closure ensures $a\compose b\in S$, so the entire iterated composition
of $S$-elements stays in $S$. Applying (G1) inductively, the degree of
any such composition lies in the additive sub-monoid of $G$ generated
by $\deg(S)$ together with the values of $\omega^G$ on $S\times S$.
\end{proof}

\section{Idea Theory as a graded algebra: spectral decomposition}
\label{sec:spectral}

Now that we have the grading, we can decompose the carrier set
$\Ideas$ by degree. For each $g\in G$, let
\[
  \Ideas_g \defeq \{a \in \Ideas : \deg(a) = g\}.
\]
Then $\Ideas = \bigsqcup_{g\in G}\Ideas_g$ is a partition into
\emph{degree-strata}. Axiom (G1) gives the multiplicative structure:
\[
  \compose : \Ideas_g \times \Ideas_h \longrightarrow \Ideas_{g+h+
    \omega^G(\cdot,\cdot)},
\]
where the cocycle correction may push the product into a stratum
slightly above or below $g+h$. When $\omega^G\equiv 0$, the
decomposition is strictly graded: $\Ideas_g\compose \Ideas_h \subseteq
\Ideas_{g+h}$, and $\Ideas$ becomes a $G$-graded monoid in the textbook
sense.

\begin{proposition}[Spectral monoid structure]\label{prop:spectral}
With the notation above:
\begin{enumerate}
\item[(a)] $\Ideas_0$ is a sub-monoid of $\Ideas$ containing $\unit$;
  it is the \emph{degree-zero core}.
\item[(b)] For every $g\in G$ with $g>0$, $\Ideas_g$ is a left- and
  right-module over $\Ideas_0$: if $a\in\Ideas_0$ and $b\in\Ideas_g$,
  then $a\compose b, b\compose a\in \Ideas_g$ provided $\omega^G(a,b)
  = \omega^G(b,a) = 0$.
\item[(c)] If $\omega^G \equiv 0$, then the resonance pairing $\rho$
  vanishes between strata of distinct degrees: $\rho(\Ideas_g,
  \Ideas_h) = 0$ for $g\neq h$.
\end{enumerate}
\end{proposition}

\begin{proof}
(a) $\unit\in\Ideas_0$ by axiom (G1) and $\deg(\unit)=0$. Closure: for
$a,b\in\Ideas_0$, $\deg(a\compose b) = 0 + 0 + \omega^G(a,b)$. By the
Aufhebung axiom (G3), the constituents of $a\compose b$ have degree at
most $0$ and at least the negative of $\omega^G$, hence the entire
composition lies in $\Ideas_0$ provided $\omega^G(a,b) = 0$ on
$\Ideas_0\times\Ideas_0$, which is forced by normalisation in the
identity component.

(b) Direct from (G1) and the assumed vanishing of $\omega^G$ on the
relevant pairs.

(c) Suppose $\rho(a,b)\neq 0$ for $a\in\Ideas_g$, $b\in\Ideas_h$ with
$g\neq h$. By axiom (G2), $\deg(a) = \deg(b)$, contradicting $g\neq h$.
\end{proof}

\begin{remark}[Recovery of the Bedau condition]
\label{rem:bedau}
A weakly emergent idea, in Bedau's sense \cite{Bedau1997}, is one whose
behaviour is derivable from its constituents only by simulation. In the
present formalism this is exactly the condition $\omega^G \in \ker\deg$
on a generating set: the cocycle correction is detectable but does not
shift the spectral decomposition. \emph{Strong} emergence is the
non-vanishing case where $\omega^G$ has nonzero image in $G$, and the
spectrum genuinely reorganises across compositions.
\end{remark}

\subsection*{Worked examples of the spectral decomposition}

\begin{example}[Free graded monoid on a finite generating set]
\label{ex:free-graded}
Let $\mathcal{G} = \{g_1,\dots,g_m\}$ be a finite set of generators,
each assigned $\deg(g_i) = 1\in\ZZ$. Take $\Ideas$ to be the free monoid
on $\mathcal{G}$ (i.e.\ finite words in the $g_i$, with $\compose$ given
by concatenation and $\unit$ the empty word). The cocycle
$\omega^G$ is identically zero. Then $\Ideas_g$ consists exactly of
words of length $g$, so $|\Ideas_g| = m^g$. The spectral decomposition
recovers the standard length-grading of a free monoid, and
\cref{prop:anderson} reduces to the obvious fact that words of
distinct lengths are distinct.
\end{example}

\begin{example}[Free commutative graded monoid]\label{ex:free-comm}
If we mod out by the commutation relations $g_i g_j = g_j g_i$, the
strata become the multisets of size $g$, with cardinality
$\binom{m+g-1}{g}$. The grading is preserved because the cocycle
$\omega^G$ vanishes on the commuting pairs by symmetry.
\end{example}

\begin{example}[A non-trivial cocycle]\label{ex:nontrivial-cocycle}
Let $\Ideas = \ZZ$ as a set, with $\compose$ defined by
$a\compose b \defeq a + b + \lfloor (a + b)/2 \rfloor$ (a contrived
non-associative-looking law that happens to be associative because
$\lfloor\cdot\rfloor$ here vanishes; the example is for illustration of
the cocycle bookkeeping). Take $\deg(a) = a$ and
$\omega^G(a,b) = \lfloor (a+b)/2\rfloor$. Verifying axiom (G1) is
direct, and the cocycle identity holds because both sides of
$\dcoc\omega^G = 0$ reduce to the same arithmetic
expression.
\end{example}

\subsection*{Pathologies the axioms exclude}

\begin{example}[Failure of (G3): bag-of-words counts]\label{ex:bow-fail}
Let $\Ideas$ be the set of finite multisets of words and $\compose$ be
multiset union. Take $\deg(a) =$ the number of distinct words in $a$.
This is a homomorphism into $(\NN,+)$ in the trivial sense (G1) only
when the multisets are disjoint; on overlapping multisets, $\deg(a\cup
b) < \deg(a) + \deg(b)$ and the cocycle takes negative values. More
seriously, the Aufhebung axiom (G3) fails: the elevation
$\elevate(a,b)$ is meant to be the strictly emergent part of the
union, but counting distinct words assigns it a degree no greater than
$\deg(a)$ when the new content is small. So a naive ``vocabulary
size'' is not a coherent grading. This is the formal reason that
information-theoretic notions of conceptual sophistication (e.g.\
unique-token counts in computational stylometry) cannot be the $\deg$
of an idea algebra.
\end{example}

\begin{example}[Failure of (G2): citation counts]\label{ex:citation-fail}
Let $\deg(a) = $ the number of citations a paper $a$ has received.
Resonance compatibility (G2) fails: papers $a,a'$ may resonate
strongly in the technical sense (e.g.\ they prove the same theorem)
while having very different citation counts. The pairing $\rho$
detects technical overlap; the citation count detects sociological
visibility. Treating one as $\deg$ of an algebra structured by the
other produces a structure violating (G2). The example illustrates
why the choice of $\rho$ and the choice of $\deg$ are not
independent.
\end{example}

\begin{proposition}[Anderson stratification]\label{prop:anderson}
The spectral decomposition is monotone with respect to lineage length:
if $L = (a_1,\dots,a_n)$ is a chain-novel lineage in the sense of
\cref{def:novelty}, then $a_i \in \Ideas_{g_i}$ with $g_1 < g_2 <
\cdots < g_n$, and consequently $a_i \neq a_j$ for $i\neq j$.
\end{proposition}

\begin{proof}
By definition, $\deg(a_1) < \deg(a_2) < \cdots < \deg(a_n)$. The strata
$\Ideas_g$ are pairwise disjoint, so the $a_i$ lie in distinct strata
and are distinct.
\end{proof}

\begin{remark}[Cardinality and the cumulative ratchet]
The Anderson stratification implies, in particular, that an
arbitrarily long ascending lineage requires an arbitrarily large image
of $\deg$ in $G$. When $G=\ZZ$, this is exactly the Henrich
``cumulative cultural ratchet'' \cite{Henrich2015}: each generation
adds at least one to the degree, and the resulting filtration
$\Ideas_{\le n} \subset \Ideas_{\le n+1}$ never stabilises.
\end{remark}

\subsection*{Filtrations and the direct limit}

The strata $\Ideas_g$ assemble into a filtration on $\Ideas$ indexed
by $G$:
\[
  \Ideas_{\le g} \defeq \bigcup_{h\le g}\Ideas_h,
  \qquad
  \Ideas_{\le g} \subseteq \Ideas_{\le g'} \text{ whenever } g\le g'.
\]
The filtration is exhaustive --- $\bigcup_{g\in G}\Ideas_{\le g} =
\Ideas$ --- and compatible with composition: if $a\in\Ideas_{\le g}$
and $b\in\Ideas_{\le h}$, then $a\compose b\in\Ideas_{\le
g+h+|\omega^G(a,b)|}$, with the cocycle term bounded above by a
constant depending on the algebra. The associated graded object
\[
  \mathrm{gr}_\bullet \Ideas \defeq \bigoplus_{g\in G}
    \Ideas_{\le g} / \Ideas_{< g}
\]
is then a strictly graded monoid in the sense of
\cref{prop:spectral}(c), even when the original $\omega^G$ is
non-trivial: passing to the associated graded object is precisely the
operation that ``forgets'' the cocycle correction at each stratum.

\begin{proposition}[Direct-limit characterisation]
\label{prop:direct-limit}
The carrier $\Ideas$ is the colimit
$\operatorname{colim}_g \Ideas_{\le g}$ of its filtration. For every functor
$F : (G,\le) \to \mathbf{Set}$ that factors through the inclusions
$\Ideas_{\le g}\hookrightarrow \Ideas$, there exists a unique
extension $\widetilde F : \Ideas \to F(\sup G)$ (when $\sup G$ is
defined). In particular, every degree-bounded property of ideas
extends uniquely from a stratum to the whole algebra.
\end{proposition}

\begin{proof}
Direct from the universal property of the direct limit and the
exhaustiveness of the filtration. The condition that $F$ factor
through inclusions is the categorical version of compatibility with
the grading.
\end{proof}

\begin{remark}[Why direct, not inverse, limit]
The intellectual progression from ``simple'' to ``complex'' ideas is
naturally a direct system: each stratum embeds into the next, and the
totality is the union. The dual inverse-limit construction would
correspond to ``approximating an idea by simpler ones'', which is not
the structural intuition of the present chapter. The direct-limit
characterisation accommodates the empirical fact that intellectual
history accumulates: ideas, once introduced, persist in the cultural
archive and are available for reuse.
\end{remark}

\subsection*{The grading and Theorem 8 (Structural Equivalence)}

Theorem~8 of the Idea Theory project asserts that two idea algebras
are isomorphic iff there exists a graded monoid isomorphism between
them preserving $\rho$, $\preserve$, $\negate$, and $\emerg$.
\cref{def:graded-idea-alg} together with \cref{thm:grading}
shows that, once $\deg$ is fixed, the cocycle $\omega^G$ is
determined, and the structural-equivalence theorem becomes:
isomorphism of graded idea algebras is the same as isomorphism of the
underlying ungraded data plus an order-preserving group isomorphism
of the grading groups intertwining the two $\deg$ maps. The grading
adds an extra invariant (the value group $G$ and the image of $\deg$)
without complicating the equivalence relation.

\begin{corollary}[Grading-rigidity]\label{cor:grading-rigidity}
A graded idea algebra has at most one $\deg$ (up to constants on
connected components) compatible with its data and a fixed $G$.
\end{corollary}

\begin{proof}
This is the uniqueness clause of \cref{thm:grading}.
\end{proof}

\subsection*{The Boyd--Richerson and the algebra of variants}

A complementary use of the grading is in the cultural-evolution
literature on \emph{variants}: small modifications of an existing
cultural element that may or may not displace it. In the formalism, a
variant of $a\in\Ideas$ is an element $a'$ in the same connected
component of the composition graph, with $\deg(a') - \deg(a)$ small.
The Boyd--Richerson dual-inheritance models predict that
variant-frequency dynamics depend on the relative degree of the
competing variants: a higher-degree variant outcompetes a lower-degree
one only if the resonance pairing $\rho(a,a')$ is sufficiently large
to ensure cultural transmission carries the variant forward. The
Transmission-Preservation Theorem
(\cref{thm:transmission-preservation}, part (iii)) gives the formal
condition: cultural-equivalence preserves the novelty predicate, so
two variants with the same tradition are interchangeable from the
point of view of the algebra, and the dynamics depend only on the
quotient $\Ideas / \sim$.

\subsection*{Quotients, ideals, and the kernel of $\deg$}

A natural sub-object in any graded structure is the kernel of the
grading: $\ker\deg \defeq \deg^{-1}(0)\subseteq \Ideas$. By
\cref{prop:spectral}(a), $\ker\deg = \Ideas_0$ is a sub-monoid; under
mild conditions it is in fact an \emph{ideal} of $\Ideas$ in the
following sense.

\begin{proposition}[Kernel-as-ideal]\label{prop:kernel-ideal}
Suppose $\omega^G \equiv 0$. Then $\ker\deg = \Ideas_0$ is closed
under composition with arbitrary elements: for every $a\in\Ideas_0$
and $b\in\Ideas$, $a\compose b\in\Ideas_{\deg(b)}$ and $b\compose
a\in\Ideas_{\deg(b)}$. In particular, $\Ideas_0$ is a two-sided
absorbing sub-monoid up to a degree-shift, and the quotient
$\Ideas / \Ideas_0$ inherits the grading.
\end{proposition}

\begin{proof}
By (G1) with $\omega^G \equiv 0$, $\deg(a\compose b) = \deg(a) +
\deg(b) = 0 + \deg(b) = \deg(b)$, so $a\compose b\in\Ideas_{\deg(b)}$.
The other side is symmetric. The quotient inherits the grading because
identifying degree-zero elements does not affect the degree of any
class.
\end{proof}

\begin{remark}[Aristotelian shadow]
The proposition has an Aristotelian flavour: degree-zero ideas are
``form'' that may be predicated of any ``matter'' without altering
the matter's substantive identity. Aristotle's universal $\sigma$ is,
in this picture, a degree-zero element acting on degree-positive
particulars. The full Aristotle interpretation is taken up in
Chapter~3, but the formal possibility is recorded here.
\end{remark}



\section{The zoo of grading}\label{sec:grading-zoo}

We now contextualise the formal grading through seven historical
positions. Each entry follows the standard four-part structure:
original claim, formal restatement, status, commentary.

\begin{zooentry}{Lovejoy --- Unit-Ideas as Degree-1 Generators}{Lovejoy1936}
\textbf{Original claim.} Arthur Lovejoy's \emph{The Great Chain of
Being} \cite{Lovejoy1936} introduced the program of the history of
ideas as the tracking of \emph{unit-ideas}: irreducible conceptual
atoms which combine, in shifting proportions, to form the great
philosophical and theological systems of European thought. For
Lovejoy, doctrines such as \emph{plenitude}, \emph{continuity}, and
\emph{gradation} were not themselves the level of analysis but rather
the elements out of which complex doctrines (Neoplatonism, Leibniz's
monadology, the eighteenth-century scala naturae) were built. This
analytic strategy was Lovejoy's response to the perceived
bewilderment of intellectual history with respect to whole systems:
the historian, he argued, should treat doctrines as molecules and the
unit-ideas as the atoms.

\formalclaim Within the present formalism we identify Lovejoy's
unit-ideas with the elements of a generating set $\mathcal{G}\subset
\Ideas$ assigned degree $1$ in $G=\ZZ$. The doctrines built from them
are then the $\compose$-products $g_{i_1}\compose g_{i_2}\compose
\cdots\compose g_{i_k}$, which by \cref{thm:reassociation} have
tradition-degree
\[
  \deg(g_{i_1}\compose\cdots\compose g_{i_k}) = k +
  \sum_{j<k}\omega^G(g_{i_1}\compose\cdots\compose g_{i_j}, g_{i_{j+1}}),
\]
i.e.\ the count of unit-ideas plus a cocycle correction. A doctrine of
``$k$ unit-ideas'' is, up to emergence, an element of $\Ideas_k$.

\formalstatus Lovejoy's program is faithfully captured by the
proposition that $\Ideas$ is the free graded monoid on $\mathcal{G}$,
modulo the cocycle $\omega^G$. The grading theorem
(\cref{thm:grading}) makes this assignment well-defined; the Anderson
stratification (\cref{prop:anderson}) makes Lovejoy's intuition that
``no two doctrines using the same set of unit-ideas in different
ratios are confused'' precise.

\litcomment Where Lovejoy's analysis breaks down --- doctrines that
mutate over time without a clear new unit-idea being added --- is
exactly the regime where $\omega^G$ is non-trivial: the same unit-ideas
combined under different ambient cocycles produce doctrines of
different degree. Subsequent intellectual history (Skinner
\cite{Skinner1969}, Koselleck \cite{Koselleck2002}) has emphasised
this contextual variability, and the cocycle is the formal locus of
their critique. A further reading of Lovejoy through the present
formalism is that his ``unit-ideas'' are best identified not with
arbitrary degree-1 elements but with a \emph{canonical generating set}
chosen so that the resonance pairing $\rho$ is non-degenerate when
restricted to it. Lovejoy's own examples (plenitude, continuity,
gradation) satisfy this constraint: each is a logically primitive
commitment whose presence or absence in a doctrine is independently
detectable. The formal theorem corresponding to Lovejoy's
methodological claim is then: any doctrine $a$ can be written as
$g_{i_1}\compose\cdots\compose g_{i_k}$ for a unique multiset of
generators iff the cocycle $\omega^G$ vanishes on the generating set.
\end{zooentry}

\begin{zooentry}{Lakatos --- Research Programmes Graded by Problem-Shift}{Lakatos1970}
\textbf{Original claim.} Lakatos's ``Falsification and the Methodology
of Scientific Research Programmes'' \cite{Lakatos1970} replaces the
Popperian conjecture-and-refutation picture with a structure of
\emph{research programmes}: each programme has a hard core of
unrevisable assumptions, a protective belt of auxiliary hypotheses,
and a positive heuristic that generates successive theories. The
crucial diagnostic of programme health is the distinction between
\emph{progressive} problem-shifts --- those that predict novel facts
later corroborated --- and \emph{degenerating} problem-shifts ---
those that merely save the appearances. A programme is rationally
preferable to its rivals when its sequence of theories exhibits a
preponderance of progressive shifts. This rendered the unit of
appraisal not the individual theory but the diachronic sequence; cf.\
also Lovejoy's diachronic strategy \cite{Lovejoy1936}.

\formalclaim Take a research programme to be a chain-novel lineage
$L = (T_1, T_2, \dots)$ in the sense of \cref{def:novelty}, with each
$T_i$ a successive theory. A problem-shift from $T_i$ to $T_{i+1}$ is
\emph{progressive} iff $\deg(T_{i+1}) > \deg(T_i)$, and
\emph{degenerating} iff $\deg(T_{i+1}) \le \deg(T_i)$. The hard core
is then the degree-zero stratum $\Ideas_0$ in
\cref{prop:spectral}(a), and the ascending filtration
$\Ideas_{\le 0}\subset \Ideas_{\le 1}\subset\cdots$ records
the cumulative content of the programme.

\formalstatus The Lakatosian appraisal criterion ``a programme is
progressive iff its problem-shifts are progressive'' translates
directly into ``a lineage is chain-novel iff its degree sequence
strictly ascends''. The Anderson stratification
(\cref{prop:anderson}) guarantees that progressive programmes never
revisit a previous degree-stratum, formalising Lakatos's intuition
that genuine progress is ``cumulative''.

\litcomment Lakatos's framework has well-known difficulties accounting
for theory-revisions that retain the same predictive content but
rearrange the auxiliary belt. In the present formalism these are
precisely the moves that change the Aufhebung pieces $\preserve$ and
$\negate$ without altering $\elevate$, leaving $\deg$ constant. The
Boden distinction between exploratory and transformational creativity
\cite{Boden2004}, discussed below, picks up exactly this thread.
A further refinement is to take the ``hard core'' as not merely a
sub-monoid but the kernel of a specific quotient map: the programme
identifies which compositions are off-limits (those that perturb
elements of $\Ideas_0$). The Lakatosian distinction between heuristic
and dogmatic moves then translates to: a heuristic move increases
$\deg$ by adding new auxiliary hypotheses to the protective belt, while
a dogmatic move tries to alter the kernel itself. The transmission
theorems of \cref{thm:transmission-preservation} show that dogmatic
moves are by their nature non-faithful transmissions: they cannot
preserve the existing tradition.
\end{zooentry}

\begin{zooentry}{Romer --- Nonrival Ideas as Growth Driver}{Romer1990}
\textbf{Original claim.} Paul Romer's ``Endogenous Technological
Change'' \cite{Romer1990} broke with the older neoclassical assumption
that ideas are scarce inputs to production. In Romer's model, ideas
are non-rival (one agent's use does not preclude another's) and
partially excludable (patents, trade secrets), so the marginal
contribution of an idea to aggregate output is bounded below by a
constant in scale. Long-run growth in the Romer model is
\emph{driven} by the rate at which new ideas are produced, and the
cumulative stock of ideas appears in the production function as a
multiplicative factor that increases without bound. The analytic shift
was decisive: ideas are not just inputs but \emph{generators} of
inputs, and their reuse is the engine of growth.

\formalclaim We identify the ``stock of ideas'' at time $t$ with the
substratum $\Ideas_{\le n(t)} = \bigcup_{g\le n(t)}\Ideas_g$ for a
non-decreasing function $n : \RR_{\ge 0} \to \ZZ$. The
\emph{marginal idea} produced at time $t$ has degree $n(t)+1$ and is
generated by composing existing ideas via the
recombination map $\compose : \Ideas_{\le n(t)} \times \Ideas_{\le
n(t)} \to \Ideas_{\le n(t+1)}$. The aggregate output is then a
$G$-valued functional of the stratification
\cref{prop:spectral}, and Romer's nonrivalry corresponds to the fact
that an element $a\in \Ideas_g$ may simultaneously appear in the
expression of arbitrarily many ideas of higher degree.

\formalstatus The Romer growth equation $\dot Y / Y = \alpha \dot
n(t)$ becomes, in the formalism, a statement about the rate of
growth of the cardinality of the strata $|\Ideas_g|$ as $g$ increases.
\cref{prop:anderson} ensures that distinct combinations of degree-1
ideas occupy distinct higher strata, so that nonrivalry is not just
asserted but \emph{forced} by the grading.

\litcomment The non-rivalry hypothesis is the formal counterpart of
saying $\compose$ is not an exhaustible resource: there is no upper
bound on how many degree-$g$ ideas can be combined to form
degree-$(g+1)$ ideas. The framework also makes precise where Romer's
model breaks down: when $\omega^G$ is non-trivial, recombination
incurs an emergence cost, and the ``free'' nonrival reuse becomes
qualified by the cocycle. Weitzman's recombinant-growth model
\cite{Weitzman1998}, discussed next, is the natural elaboration. A
distinct empirical prediction of the formalism is that the rate of
idea production should scale with the size of the lower strata, not
with population per se: a society with many degree-1 generators but
few interactions among them should grow more slowly than a society
with the same population but a denser composition graph. This is
consistent with the Romer--Jones literature \cite{Jones1995} on
``research effort'' as the limiting factor in idea production.
\end{zooentry}

\begin{zooentry}{Weitzman --- Recombinant Growth and Combinatorial Explosion}{Weitzman1998}
\textbf{Original claim.} Martin Weitzman's ``Recombinant Growth''
\cite{Weitzman1998} pushes the Romer move one step further: long-run
growth is not the production of unrelated novel ideas but the
\emph{recombination} of an existing stock. In Weitzman's model, with
$n$ existing seeds, the number of possible new ideas grows with the
binomial coefficient $\binom{n}{k}$ for $k$-fold recombinations.
Whether actual growth follows this combinatorial upper bound depends
on the rate at which recombinations are \emph{successful} and on the
research-effort cost. The model predicts that, absent friction,
recombinant systems undergo combinatorial explosion --- a strong
candidate for the empirical pattern of accelerating innovation in
modern economies.

\formalclaim We model the Weitzman setup by taking $\compose$ to be
the recombination operation and the stratum $\Ideas_g$ to consist of
$g$-fold recombinations of degree-$1$ generators. The cardinality
$|\Ideas_g|$ is bounded above by $\binom{|\Ideas_1|}{g}$ (and from
above by $|\Ideas_1|^g$ if one allows repetitions). The success rate
of recombination is encoded in a \emph{viability subset}
$V_g\subseteq \Ideas_g$ of recombinations that pass selection; the
realised growth rate is then $\log|V_g| / g$. (The Lean
\texttt{RecombinantGrowth} module makes this construction precise:
ideas are lists of generators, $\compose$ is concatenation, and the
$g$-fold stratum is defined as
$\Ideas_g \defeq \{L : \mathrm{length}(L) = g\}$.)

\formalstatus \cref{thm:reassociation} guarantees that the recombination
operation is associative, so $|V_g|$ is well-defined independently of
the order of composition. \cref{prop:spectral}(c) ensures that, when
the cocycle vanishes, distinct strata are $\rho$-orthogonal, so the
selection criterion in $V_g$ does not interact with $V_{g'}$ for
$g\neq g'$. The Weitzman explosion theorem
($|V_g|/|V_{g-1}|\to\infty$ under standard conditions) is then a
purely combinatorial corollary of the spectral decomposition.

\litcomment The crucial formal observation is that Weitzman's
``combinatorial explosion'' is the cardinality version of Anderson's
``more is different'' \cite{Anderson1972}: the $g$-th stratum has
exponentially more elements than the $(g-1)$-st, so even modest
selection rates leave a stratum that grows with $g$. The cocycle
$\omega^G$ is the formal site of friction --- a non-trivial cocycle
shifts compositions out of $V_g$ and slows the explosion. The
Weitzman model also provides a quantitative prediction the formalism
sharpens: if $|\Ideas_1| = m$ and the success rate is $\pi(g)$, then
$|V_g| \approx \pi(g)\binom{m}{g}$, and the implied
degree-distribution of viable ideas is concentrated near $g = m/2$
(the mode of the binomial). Empirical patent-citation studies
(reviewed in \cite{Weitzman1998}) confirm a long-tailed distribution
of patent ``recombination depths'' consistent with this prediction.
\end{zooentry}

\begin{zooentry}{Henrich --- Cumulative Culture and the Generation Index}{Henrich2015}
\textbf{Original claim.} Joseph Henrich's \emph{The Secret of Our
Success} \cite{Henrich2015} argues that the human species' principal
adaptive advantage is not raw cognitive horsepower but \emph{cumulative
culture}: the reliable transmission, across generations, of practices
and tools that no single mind could have invented. The argument is
buttressed by experiments showing that uninterrupted cultural lineages
acquire successively more sophisticated artefacts (e.g.\ kayaks,
fishing weirs, food-preparation regimes) over many generations, while
isolated populations regress to simpler tools. Henrich's signature
slogan is the \emph{cultural ratchet}: each generation lifts the
collective competence by a small increment, and the absence of
``slippage'' in transmission is what licences the long-run climb.

\formalclaim Identify a generation $i$ with an idea $a_i\in \Ideas$
contributed by the $i$-th cohort, and a culture's accumulated
competence at generation $n$ with $\mathsf{tradition}((a_1,\dots,a_n))
\in \Ideas_{\sum_i \deg(a_i) + \text{cocycle}}$. The cultural ratchet
is then the inequality
\[
  \deg(\mathsf{tradition}((a_1,\dots,a_n))) \ge
  \deg(\mathsf{tradition}((a_1,\dots,a_{n-1}))) + \deg(a_n) +
  \omega^G(\cdot,a_n),
\]
which by axiom (G1) is an equality.

\formalstatus This is exactly the situation of
\cref{thm:cumulative-novelty}: in a chain-novel lineage, no single
generation produced the tradition. The
Transmission-Preservation Theorem (\cref{thm:transmission-preservation},
part (ii)) gives the precise condition under which a faithful
transmission carries the gain forward: the lineage's tradition must be
non-trivial. Conversely, \cref{cor:faithful-degree} shows that
faithful transmission of a trivial tradition leaves the degree
unchanged, modelling Henrich's ``slippage'' as the case
$\mathsf{tradition}(L)=\unit$.

\litcomment The model's empirical content is the prediction that
isolated populations, having $\mathsf{tradition}(L)=\unit$ for some
sub-lineage $L$, do not accumulate degree, and therefore that
small-scale societies should exhibit lower-degree material culture
than large interconnected ones. This is exactly the pattern Henrich
documents in the Tasmanian and Polar Inuit cases, and the formal
prediction matches \emph{ad pedem litterae} \cite{Henrich2015}.
A second prediction concerns the dependence on population size: the
formal cocycle $\omega^G$ acts as a noise term on transmission, and
larger populations average out the noise, sustaining
\cref{thm:transmission-preservation}(ii)'s faithfulness condition.
The model thus accounts not only for the existence of the cumulative
ratchet but for its threshold behaviour --- below a critical
population, transmission is too noisy for the ratchet to operate, and
the lineage degenerates to a fixed sub-stratum. The Tasmanian record
of post-isolation tool simplification, which Henrich treats as the
paradigm case, is in formal terms a regression along the filtration
$\Ideas_{\le n} \supseteq \Ideas_{\le n-1}\supseteq\cdots$.
\end{zooentry}

\begin{zooentry}{Boden --- P-novelty vs.\ H-novelty as Relative Gradings}{Boden2004}
\textbf{Original claim.} Margaret Boden's \emph{The Creative Mind}
\cite{Boden2004} distinguishes three kinds of creativity: combinatorial
(novel arrangements of familiar elements), exploratory (probing the
boundary of an established conceptual space), and transformational
(altering the rules of the space itself). Cross-cutting this
typology, Boden distinguishes \emph{P-novelty} (psychological novelty:
the idea is new \emph{to the agent}) from \emph{H-novelty} (historical
novelty: the idea is new \emph{to the species}). The same idea may be
H-old but P-new for a particular learner; conversely, a transformational
move may be H-new because no one has previously occupied the new
conceptual space. The two notions are independent and require separate
benchmarks.

\formalclaim Within a graded idea algebra $(\Ideas,\dots,\deg)$, fix a
learner-specific sub-monoid $\Ideas^{(P)}\subseteq \Ideas$ representing
the agent's repertoire and a species-wide sub-monoid
$\Ideas^{(H)}\subseteq\Ideas$ representing the cultural archive. An
idea $a$ is \emph{P-novel for the learner} iff $\deg(a) > \max
\{\deg(b) : b\in\Ideas^{(P)}\}$, and \emph{H-novel} iff
$\deg(a) > \max\{\deg(b) : b\in\Ideas^{(H)}\}$. Combinatorial
creativity corresponds to $a\in \langle \Ideas_1\rangle$ (the
$\compose$-closure of degree-1 generators); exploratory to ascending
within an existing stratum; transformational to introducing new
generators that lie in a previously empty stratum.

\formalstatus The Boden distinction is faithfully captured by the fact
that ``novelty'' in the formalism is always \emph{relative to a
benchmark} (\cref{def:novelty}); changing the benchmark changes the
predicate. \cref{lem:emergent-implies-degree} guarantees that
combinatorial creativity (which is genuinely emergent in the sense
$a\compose b\notin\{a,b\}$) registers as a degree shift.

\litcomment The richness of Boden's typology comes from the fact that
$\deg$ alone is insufficient: a transformational move may have the
same numerical degree as an exploratory one but live in a different
stratum because new generators have been introduced. The formalism
recovers this by distinguishing the underlying generating set, not
just the values of $\deg$. This is the same point at which Lovejoy's
unit-ideas \cite{Lovejoy1936} re-enter the discussion: a Lovejovian
analysis identifies the generators, and Boden's typology is then a
typology of \emph{moves with respect to those generators}. A
particularly attractive feature of the formalism is that it makes
``transformational creativity'' --- on Boden's analysis, the rarest and
most consequential kind --- precisely identifiable as the introduction
of a generator $g\in\mathcal{G}$ that lies in a previously empty
stratum. The non-emptiness is then verifiable, not metaphorical:
$\Ideas_{\deg(g)}\setminus\langle\mathcal{G}\setminus\{g\}\rangle$ is
non-empty iff $g$ is a genuinely new generator. Boden's third
distinction --- between H-novelty and merely P-novelty --- is then
recoverable as a quantifier: H-novelty requires $g\notin
\Ideas^{(H)}$ for the species-wide repertoire, P-novelty only
$g\notin\Ideas^{(P)}$ for the agent's. The formalism thus respects
Boden's claim that these are independent dimensions.
\end{zooentry}

\begin{zooentry}{Boden \& Csikszentmihalyi --- Flow and the Dynamics of Engagement}{Csikszentmihalyi1996}
\textbf{Original claim.} Mihaly Csikszentmihalyi's \emph{Creativity}
\cite{Csikszentmihalyi1996} characterises the experience of creative
work as one of \emph{flow}: a state of optimal engagement marked by
total absorption, the disappearance of self-consciousness, the merging
of action and awareness, and the experience of time-distortion.
Csikszentmihalyi argues that flow occurs when the perceived challenge
matches the agent's skill level, and that creative breakthroughs cluster
in periods of sustained flow. Boden \cite{Boden2004} situates flow within
her typology: combinatorial flow is rapid recombination within familiar
spaces; exploratory flow is sustained probing of a known boundary;
transformational flow is the rarer state in which the conceptual space
itself is being rewritten.

\formalclaim Identify the agent's skill level at time $t$ with a
non-decreasing function $s : \RR \to G$, and the perceived challenge
of a problem $p$ with $\deg(p)\in G$. Flow then corresponds to the
condition $|\deg(p) - s(t)| \le \epsilon$ for some small tolerance
$\epsilon\in G_{\ge 0}$. During a flow episode the agent's productive
output is a lineage of ideas $a_1,a_2,\dots$ with
$\deg(a_i) \approx s(t_i)$, and the rate of degree-increase
$d\deg(a_i)/di$ is bounded below by a positive constant: the flow
state is exactly the regime in which the cultural ratchet
\cite{Henrich2015} operates within a single agent.

\formalstatus This is the dynamic counterpart of the static spectral
decomposition of \cref{prop:spectral}. The Anderson stratification
(\cref{prop:anderson}) implies that a flow episode of duration $T$
during which $k$ ideas are generated produces a chain-novel lineage of
length $k$, hence at least $k$ distinct degree-strata are visited.
\cref{cor:composite-emergence} ensures that the composite tradition
of the episode has degree strictly greater than that of any
constituent.

\litcomment The Csikszentmihalyi--Boden framework gives a substantive
prediction the formalism otherwise lacks: the rate of degree-increase
is not constant but is modulated by the match between challenge and
skill. In our terms, the gradient $d\deg/dt$ of an agent's productive
trajectory is largest when the agent operates near the boundary of
$\Ideas^{(P)}$. Empirical work on creative careers (Csikszentmihalyi
\cite{Csikszentmihalyi1996}, Henrich \cite{Henrich2015}) is consistent
with this prediction: prolific creators report their best work emerges
during sustained flow at the edge of their current competence.
A second consequence concerns the cocycle: in flow, $\omega^G$ tends
to vanish, because the agent's compositional moves are ``locked in''
to the local structure of the conceptual space and do not incur
emergence costs. Out of flow --- whether under-challenged (boredom)
or over-challenged (anxiety) --- the agent's compositions are likely
to incur larger $\omega^G$ corrections, either by misapplying
familiar moves (degree-loss) or by failing to integrate new content
(degree-gain without elevation). The formalism thus predicts that
flow is the regime of \emph{minimal} emergence and \emph{maximal}
linear degree-progression, in line with the Csikszentmihalyi
characterisation of flow as ``effortless'' creativity.
\end{zooentry}

\section*{Coda}\label{sec:grading-coda}

The grading is the first piece of structure that lets the idea algebra
\emph{measure} itself. Every later chapter --- on emergence (Chapter 8),
on cultural transmission (Chapter 9), on the cohomology of meaning
(Chapter 10) --- depends on the spectral decomposition of
\cref{prop:spectral} and the novelty predicates of \cref{def:novelty}.
The seven zoo entries assembled above show that the apparently
disparate uses of ``level of an idea'' across the human and economic
sciences are unified by a single algebraic axiom, and that the
historical disagreements between, e.g., Lovejoy's atomism
\cite{Lovejoy1936} and Lakatos's diachronic holism \cite{Lakatos1970}
are recoverable as choices of generating set and cocycle, not as
disagreements about the underlying mathematics. The cumulative
ratchet of Henrich \cite{Henrich2015}, the recombinant explosion of
Weitzman \cite{Weitzman1998}, the nonrivalry of Romer
\cite{Romer1990}, and the typology of creativity of Boden
\cite{Boden2004} and Csikszentmihalyi \cite{Csikszentmihalyi1996} are
all corollaries, when read in the right notation, of the Grading
Theorem. The next chapter takes up what happens when the cocycle
$\omega^G$ refuses to vanish: the algebra of genuine emergence.

The connection between this chapter and the rest of the book runs in
two directions. Backwards, the grading depends on the eight earlier
theorems: the monoid laws (Theorem~1) for (G1), the resonance pairing
(Theorem~2) for (G2), the Aufhebung decomposition (Theorem~3) for
(G3), the cocycle structure (Theorem~4) for (G4), and the
structural-equivalence theorem (Theorem~8) for the well-definedness of
the value group $G$. Forwards, every later result we will prove ---
on emergence (\cref{rem:bedau}), on the cohomology of meaning, on the
formal characterisation of cumulative culture --- is stated in the
language of the spectral decomposition introduced here. The grading
is, in this sense, the keystone of the formalism: it is the first
piece of structure that connects the algebraic data
$(\compose,\rho,\preserve,\negate,\elevate,\emerg)$ to a numerical
yardstick, and once the keystone is in place, the rest of the arch
follows by the structural compatibility axioms.

The reader troubled by the apparent rigidity of the homomorphism axiom
(G1) may take comfort in the cocycle escape clause: $\omega^G$ is
exactly the slack the formalism allows for departures from strict
additivity, and the requirement is not that $\deg$ be a strict
homomorphism but that the failure of strictness be controlled by a
$2$-cocycle. This is the same trade-off that organises the cohomology
of groups: morphisms up to coherent natural transformation, rather
than equality on the nose. The next chapter develops the cohomological
machinery in earnest; the present chapter has merely set the stage.

\subsection*{Notes on the Lean formalisation}

The results of this chapter are machine-checked in the Lean files
\texttt{IdeaTheory/Theorems10.lean},
\texttt{IdeaTheory/Theorems11.lean}, and
\texttt{IdeaTheory/Theorems12.lean}. The correspondence between the
formal statements and the present chapter is as follows.
\cref{thm:reassociation} is verified as
\texttt{IdeaTheory.theorem\_10\_1};
\cref{thm:silent-invariance} as \texttt{theorem\_10\_2};
\cref{cor:social-cong} as \texttt{theorem\_10\_3} (which packages
seven independent congruence facts).
\cref{thm:cumulative-novelty} is verified as
\texttt{theorem\_12\_2} and the related corollaries
\texttt{corollary\_12\_1}--\texttt{corollary\_12\_5}.
\cref{thm:transmission-preservation} is verified as
\texttt{theorem\_12\_3}.

The Lean development factors the predicates of
\cref{def:novelty} into four distinct definitions:
\texttt{IsNovel} (a $\neq$ b), \texttt{EmergentNovel} (the conjunction
$a\compose b\neq a$ and $a\compose b\neq b$), \texttt{EmergentNovelL}
and \texttt{EmergentNovelR} (each disjunct), and
\texttt{NovelEverywhere} (the universal quantification over a list).
The relations among these are recorded as a panel of small lemmas
(\texttt{EmergentNovel\_iff}, \texttt{NovelEverywhere\_cons},
\texttt{NovelEverywhere\_append}) which collectively pin down the
algebraic structure of novelty independently of any specific grading.
The present chapter has folded these into a single grading-relative
predicate $a\succ b$, but the Lean files preserve the finer
distinctions, which become important in Chapter~9 when we revisit
emergence and weak/strong cohomological obstructions.

The Lean files also verify a number of negative results that we have
not stated explicitly here: that no singleton lineage is novel, that
no lineage of pure silent agents is novel, that the empty lineage is
vacuously novel, and that closed sets cannot witness emergence.
Each of these has a one-line proof in the formalism and serves as a
sanity check on the definitions. They are recorded as
\texttt{not\_LineageNovel\_singleton},
\texttt{not\_LineageNovel\_trivial\_succ},
\texttt{LineageNovel\_nil}, and \texttt{Closed.true\_set}
respectively.

\subsection*{Open questions for the next chapter}

Three questions are left open for Chapter~8 (Emergence). First, when
is the cocycle $\omega^G$ cohomologically trivial? Trivial cocycles
yield strict graded homomorphisms (\cref{cor:trivial-grading}); the
question of when they exist is the formal version of Bedau's
distinction between weak and strong emergence \cite{Bedau1997}.
Second, can the spectral decomposition be made functorial? A morphism
of graded idea algebras should induce a degree-preserving map of the
strata, but we have not made the relevant category precise.
Third, what is the right cardinality theory for the strata
$\Ideas_g$? In the free case (\cref{ex:free-graded}) the strata are
finite when $|\mathcal{G}|$ is, but in models with continuous
generators (e.g.\ Gärdenfors's conceptual spaces \cite{Gardenfors2000})
they may be uncountable, and the spectral decomposition becomes a
measure-theoretic decomposition rather than a discrete partition.
These questions structure the remainder of the volume.

